\documentclass{letter}

\usepackage[utf8]{inputenc}
\usepackage{xcolor}
\usepackage{url}

\newcommand{\todo}[1]{\textcolor{red}{TODO: #1}}

\signature{Shing H. Zhan and Jerome Kelleher}

\address{Big Data Institute\\University of Oxford\\UK}
\begin{document}

\begin{letter}{Nature Genetics}

\opening{Dear Dr Li,}

We thank you and the reviewers for the careful and constructive evaluation of
our Article, ``A Pandemic-Scale Ancestral Recombination Graph for
SARS-CoV-2.''

\todo{Summary of the main changes.}

We have included a copy of the manuscript in which all changes have been
highlighted.

Below we respond point by point.

\medskip
\hrule
\medskip

\noindent \textbf{Response to the editor}

\noindent \textbf{Editor:} ``To guide the scope of the revisions, the editors
discuss the referee reports in detail within the team with a view to
identifying key priorities that should be addressed in revision. In this case,
we think both referees have provided constructive reviews aimed at
strengthening the analyses and improving the presentation. We particularly ask
that you demonstrate the practical utility of this method by expanding its
application to a broader range of plausible pandemic-causing pathogens, and
address all referee comments as thoroughly as possible with appropriate
revisions.''

\todo{Response.}

\noindent \textbf{Editor:} ``When re-submitting your manuscript, please ensure
that any supplementary figures and tables that are crucial to the manuscript's
conclusions are converted into Extended Data figures and tables to increase
visibility of these data. [\ldots] A maximum of ten Extended Data display items
(figures and tables) is permitted.''

\todo{Response.}

\medskip
\hrule
\medskip

\noindent \textbf{Response to Reviewer \#1}

\noindent \textbf{Reviewer \#1:} ``This manuscript introduces a method for
constructing an ARG from millions of viral genome sequences sampled during the
pandemic. They show that a `real-time' ARG can be constructed at scale and is
similarly or more accurate than existing methods for incorporating
recombination into genealogical inference. Overall, the method is well
described and the analyses are clearly documented and largely appropriate for
the purpose of evaluating the results. I have a few specific analyses that I
suggest below and consider to be essential for investigating the quality of
the resulting ARG and method.''

\todo{Acknowledgement.}

\noindent \textbf{Reviewer \#1:} ``Most critically, I wonder who this is for
exactly? As a resource, I have substantial misgivings about whether applied
public health would have any interest or practical ability to use it.
Phylogenetic trees are already challenging in many public health
applications. Furthermore, the ARG provided is not real time because there is
no included plan to update as new data are deposited. It's already far behind
the evolution of SARS-CoV-2 so won't be used soon in this regard.''

\todo{Response.}

\noindent \textbf{Reviewer \#1:} ``If this is intended as a method suitable for
genomic epidemiology applications in future pandemics, it needs to be
evaluated against a much broader range of plausible pandemic-causing pathogens
with the limitations and recommendations clearly established.''

\todo{Response.}

\medskip
\noindent \textit{Major}

\noindent \textbf{Reviewer \#1:} ``The post hoc examination of factors
consistent with increased recombination rates (case rates and viral genetic
diversity) seems reasonable, but should be much more precise. Specifically,
recombination requires that viral lineages co-circulate in space and time.
Given that most samples will have both a sampling date and a location, this
analysis should be much more granular and could be substantially more
convincing. Do the identified precise parental lineages co-circulate in region
and time? Additionally, given that recombination between closely related
samples will likely not be detectable by this and other methods, the
relationship between lineage diversity and recombination rates might be
expected as a consequence of the inference procedure. While I still think this
is useful to consider, the explanation and discussion could be more
nuanced.''

\todo{Response.}

\noindent \textbf{Reviewer \#1:} ``There has been some research in the effect
of recurrent mutli-nucleotide mutations with convincing examples at quite a
few sites (see de Maio et al.
\url{https://academic.oup.com/mbe/article/42/11/msaf272/8301027}). I would
expect that sc2ts could be impacted by these events resulting in elevated
rates of inferred recombination. This should be addressed and if possible
mitigated.''

\todo{Response.}

\noindent \textbf{Reviewer \#1:} ``The SARS-CoV-2 pandemic has largely come and
gone and global sequencing efforts are in free-fall (note that this resource
is not updated with the most recently produced data). While I appreciate the
tool could be `real time' in a new pandemic, there's no telling exactly what
the next pandemic will resemble. If this is going to be a meaningful method
going forward, I suggest this should be evaluated across simulations
encompassing varied plausible parameters. E.g., a flu pandemic is thought to
be likely, could this be used to handle the kinds of reassortment that we see
with segmented genomes?''

\todo{Response.}

\noindent \textbf{Reviewer \#1:} ``In parts of the introduction and results,
the overall tone is one that suggests this method has solved genealogical
inference for SARS-CoV-2 and to a casual reader might indicate very accurate
genealogical inference. While I do think some aspects are improved by this
work, the uncertainty is still enormous, all `validations' applied to real
data are `corroborative analyses' in reality. This work would benefit from a
substantive effort to soften and inject more circumspect phrasing
throughout.''

\todo{Response.}

\noindent \textbf{Reviewer \#1:} ``Some recommendations should be provided for
the expected practical use cases. Much of the text suggests this will be used
by researchers -- certainly I would expect population geneticists and the like
to be interested. I don't think this has much practical utility in applied
public health where it is already a struggle to incorporate single phylogenies
and instead we use broad-scale categories for analysis and reporting concerns.
If this is a new method and resource, it is critical that we establish who the
intended user-base is in reality and how they would interact with it to
address their concerns.''

\todo{Response.}

\medskip
\noindent \textit{Minor}

\noindent \textbf{Reviewer \#1:} ``The word `validated' (for example, when
considering 3seq agreement) is an overstatement. It is a comparison against an
inference method using the same underlying data. These events are
`corroborated' or something more along those lines. The phrasing `The ARG is a
comprehensive and deeply validated record of the pandemic phase' is a
overstatement for similar reasons (see also above).''

\todo{Response.}

\noindent \textbf{Reviewer \#1:} ``I am not sure if the authors will have seen
this preprint,
\url{https://www.biorxiv.org/content/10.64898/2026.02.19.706861v1.abstract},
which provides a way to visualize and explore large-scale ARGs? I realize this
was published very recently, so it may not be ready for this specific
application. One way that this resource could be made much more accessible
would be to present this ARG on the platform.''

\todo{Response.}

\noindent \textbf{Reviewer \#1:} ``One element that may be worth mentioning is
that the original RIPPLES analysis was run on the `real' UShER tree. While the
viridian tree is similar, it has not had anywhere near the same level of
community inspection and manual curation. Whether that curation impacts the
initial placements of plausible recombinants is anyone's guess, and I
certainly don't expect the authors to take on that problem, but this
distinction might be worth addressing briefly.''

\todo{Response.}

\noindent \textbf{Reviewer \#1:} ``Is there a plan to update the ARG with newly
produced data? This would certainly help establish this as a `real time'
resource.''

\todo{Response.}

\medskip
\noindent \textit{Code availability}

\noindent \textbf{Reviewer \#1:} ``The documentation is excellent. While I have
not gone through everything myself, I am confident that I could reproduce the
analyses presented in the manuscript with sufficient effort.''

\todo{Acknowledgement.}

\medskip
\hrule
\medskip

\noindent \textbf{Response to Reviewer \#2}

\noindent \textbf{Reviewer \#2:} ``Summary: In this manuscript, the authors
present both an extraordinary data resource -- an inferred Ancestral
Recombination Graph (ARG) for SARS-CoV-2 -- as well as a statistical method --
sc2ts -- for inferring ARGs using time-series sampling of pandemic-scale
genomic data. This method builds off a widely used model in population
genetics, the Li and Stephens Hidden Markov Model, that is tweaked here to
rely only on a `recombination penalty' parameter, the value of which is
determined heuristically. The inferred ARG is rigorously compared and
sanity-checked against what is known (or has been previously inferred) about
the history of the genomic evolution of SARS-CoV-2.''

\noindent ``Opinion: Although recombination between SARS-CoV-2 lineages is
rare, it has unambiguously occurred, which means that the true genetic history
of the virus is best understood as a reticulate genealogy, i.e., an ARG. An
accurate genetic history, as well as tools for inferring such histories, are
crucial not only to fundamental questions in epidemiology and evolution, but
also to the development of appropriate measures for mitigating the spread of
human pathogens. The work presented in this manuscript represents a clear step
forward as both an accessible data resource for our understanding of the
evolution of the modern diversity of SARS-CoV-2 and as a method for
understanding the evolutionary trajectory of future pandemics. The manuscript
is clear and well-written, and I had only minor suggestions for increasing the
clarity and impact of the work.''

\todo{Acknowledgement.}

\noindent \textbf{Reviewer \#2 (1):} ``Future work. The only major comments I
have deal with future work that I hope will come out of this group.''

\noindent \textbf{Reviewer \#2 (1a):} ``First, with any novel statistical
method, it is reassuring to assess its behavior on simulated datasets. In this
case, the authors are able to assess the reliability and accuracy of their
method reasonably well by simply relying on the extraordinary wealth of data
on SARS-CoV-2 evolution in lieu of simulations where the truth is known.
However, I think a future project in which the performance of sc2ts is
assessed using viral simulations would be useful, both for quantifying the
method's behavior and identifying areas of parameter space (e.g., rates of
recombination vs.\ mutation, or time-density of sampling, etc.) in which its
performance might become unreliable.''

\todo{Response.}

\noindent \textbf{Reviewer \#2 (1b):} ``Second (although following from 1a
above), it seems like much of the inference done by sc2ts relies on
heuristics. This is understandable, and indeed difficult to avoid in the
analysis of real-world empirical biological data, which is often messy or
unreliable. However, the authors note that a goal of the current work is to
prepare for the next pandemic, so I encourage them to consider publishing a
companion manuscript that walks through the decision-making process involved
in the selection of, e.g., the `k' parameter so that researchers employing
this method in the study of other pathogens might have a playbook off of which
to work.''

\todo{Response.}

\noindent \textbf{Reviewer \#2 (2):} ``Main Figure: ARGs are often difficult to
parse visually, but I feel this paper would be more compelling with a
graphical demonstration of the inferred relatedness structure enabled by
sc2ts. I wonder if the authors could adapt Figure 2 to display inferred ARG of
SARS-CoV-2? Perhaps the complexity of such a figure could be reduced by
collapsing some clades into triangles, as is often done in the phylogenetics
literature, in order to emphasize the important inferred reticulations? This
may not be possible (and it's by no means a dealbreaker for me as a reviewer),
but I think it would help demonstrate the power of this approach to a general
audience.''

\todo{Response.}

\noindent \textbf{Reviewer \#2 (3):} ``Writing. I had some minor
corrections/suggestions to increase the clarity of the work.''

\noindent \textbf{Reviewer \#2 (3a):} ``L45: unclear if `over time' means over
the course of a pandemic or over computer time in the algorithm of the
method.''

\todo{Response.}

\noindent \textbf{Reviewer \#2 (3b):} ``L161: missing word.''

\todo{Response.}

\noindent \textbf{Reviewer \#2 (3c):} ``L166: should be `855 inferred
recombination events' right?''

\todo{Response.}

\noindent \textbf{Reviewer \#2 (3d):} ``L236: it might be useful to define
`Scorpio' for the audience of a general journal like this''

\todo{Response.}

\noindent \textbf{Reviewer \#2 (3e):} ``L243: should be `inferred recombination
events' right?''

\todo{Response.}

\noindent \textbf{Reviewer \#2 (3f):} ``Throughout, it seems like sometimes the
quotations marks were facing the wrong direction (i.e., closing when they
should be opening, or vice versa), so I'd recommend the authors do a careful
proofread for this.''

\todo{Response.}

\closing{Sincerely,}

\end{letter}
\end{document}
